\section{The Parallel Tempering Algorithm}
\label{sec:algorithm}

\subsection{Notation and Setup}

Let $G = (V, E, w)$ be the census-tract adjacency graph for a state, where
$V$ is the set of tracts, $E$ is the set of shared-boundary edges, and
$w: E \to \mathbb{R}_{>0}$ is the geographic boundary-length weight.
Let $p: V \to \mathbb{Z}_{>0}$ be the 2020 Census tract populations, and let
$k$ be the target number of congressional districts.
A \emph{plan} is an assignment $\mathbf{a}: V \to \{1,\ldots,k\}$ such that
every district $d = \mathbf{a}^{-1}(j)$ is connected in $G$ and satisfies the
population balance constraint at tolerance $\varepsilon$:
\[
  \left|\frac{\sum_{v \in d} p(v)}{P/k} - 1\right| \le \varepsilon,
\]
where $P = \sum_{v \in V} p(v)$ is the state total population.
The \emph{edge cut} of a plan is
\[
  \EC(\mathbf{a}) = \sum_{\{u,v\} \in E:\; \mathbf{a}(u) \ne \mathbf{a}(v)} w(u,v).
\]

\subsection{Geometric Tolerance Ladder}
\label{sec:ladder}

\pt{} runs $N = N_{\mathrm{replicas}}$ chains indexed $i = 0, 1, \ldots, N-1$.
Replica $i$ operates at balance tolerance $\varepsilon_i$ drawn from a
\emph{geometric ladder}:
\[
  \varepsilon_i = \varepsilon_{\mathrm{cold}}
    \times \left(\frac{\varepsilon_{\mathrm{hot}}}{\varepsilon_{\mathrm{cold}}}\right)^{i/(N-1)},
  \quad i = 0, 1, \ldots, N-1.
\]
This gives $\varepsilon_0 = \varepsilon_{\mathrm{cold}}$ (the cold chain, strict statutory
tolerance) and $\varepsilon_{N-1} = \varepsilon_{\mathrm{hot}}$ (the hot chain,
relaxed tolerance).
All intermediate tolerances are strictly increasing.

\begin{definition}[Inverse tolerance / inverse temperature]
\label{def:beta}
For replica $i$ with tolerance $\varepsilon_i$, the \emph{inverse tolerance}
(the analogue of inverse temperature) is $\beta_i = 1 / \varepsilon_i$.
The cold chain has the highest $\beta$ (most constrained); the hot chain has
the lowest $\beta$ (least constrained).
\end{definition}

\paragraph{Equal swap acceptance design.}
The geometric ladder is chosen so that the expected swap acceptance rate between
adjacent replicas is approximately equal across all pairs.
For continuous Gaussian systems, equal spacing in $\ln\varepsilon$ achieves
${\approx}23\%$ swap acceptance~\citep{kone2005}.
For discrete redistricting chains, this serves as an empirical starting point.
The ladder spacing can be diagnosed from the observed swap acceptance rates in
the audit log; a rate substantially below $15\%$ suggests over-spacing
(too many replicas needed), and a rate substantially above $35\%$ suggests
under-spacing (fewer replicas would suffice).

\paragraph{Default parameters.}
\begin{center}
\renewcommand{\arraystretch}{1.25}
\begin{tabular}{@{}lll@{}}
\toprule
Parameter & Default & Notes \\
\midrule
$N_{\mathrm{replicas}}$ & 4 & 4 chains at $\varepsilon \in \{0.005, 0.013, 0.033, 0.05\}$ \\
$\varepsilon_{\mathrm{cold}}$ & 0.005 & $\pm 0.5\%$ statutory tolerance \\
$\varepsilon_{\mathrm{hot}}$ & 0.05 & $\pm 5\%$ relaxed tolerance \\
$\mathtt{swap\_interval}$ & 10 & Swap proposals every 10 steps \\
$T$ (steps) & 1000 & Total cold chain steps \\
$p$ (percentile) & 0.0 & Select minimum-EC cold chain plan \\
\bottomrule
\end{tabular}
\end{center}

\subsection{Seeding Specification}
\label{sec:seeding}

All random streams are derived from a single \texttt{base\_seed} via SHA-256
with domain-separated prefixes, following the platform-wide seed design.
Four distinct prefixes are used:

\begin{definition}[Replica Seed]
\label{def:replica-seed}
\[
  \mathtt{replica\_seed}(b, r, s) =
  \mathrm{first8}\bigl(\mathrm{SHA\text{-}256}(
    \texttt{"PT\_REPLICA\_"} \;\|\;
    \mathrm{le4}(r) \;\|\;
    \texttt{"\_"} \;\|\;
    \mathrm{le8}(s) \;\|\;
    \texttt{"\_"} \;\|\;
    \mathrm{le8}(b)
  )\bigr),
\]
where $b$ is the base seed (u64), $r$ is the replica index (u32), $s$ is the
step index (u64), $\mathrm{le4}$ and $\mathrm{le8}$ encode values as little-endian bytes,
and $\mathrm{first8}$ reads the first 8 bytes of the SHA-256 digest as a
little-endian u64.
\end{definition}

\begin{definition}[Swap Seed]
\label{def:swap-seed}
\[
  \mathtt{swap\_seed}(b, s, i) =
  \mathrm{first8}\bigl(\mathrm{SHA\text{-}256}(
    \texttt{"PT\_SWAP\_"} \;\|\;
    \mathrm{le4}(i) \;\|\;
    \texttt{"\_"} \;\|\;
    \mathrm{le8}(s) \;\|\;
    \texttt{"\_"} \;\|\;
    \mathrm{le8}(b)
  )\bigr),
\]
where $i$ is the pair index (0-based, for the swap between replicas $i$ and $i+1$).
\end{definition}

The forward and reverse RNG streams for replica $r$ at step $s$ are derived
from the replica seed:
\[
  \mathtt{rng\_fwd} = \mathrm{SmallRng::seed\_from\_u64}\!\bigl(
    \mathrm{first8}(\mathrm{SHA\text{-}256}(\texttt{"PT\_FWD\_"} \;\|\; \mathrm{le8}(r_s)))
  \bigr),
\]
\[
  \mathtt{rng\_rev} = \mathrm{SmallRng::seed\_from\_u64}\!\bigl(
    \mathrm{first8}(\mathrm{SHA\text{-}256}(\texttt{"PT\_REV\_"} \;\|\; \mathrm{le8}(r_s)))
  \bigr),
\]
where $r_s = \mathtt{replica\_seed}(b, r, s)$.

\begin{remark}
The four prefixes \texttt{"PT\_REPLICA\_"}, \texttt{"PT\_SWAP\_"},
\texttt{"PT\_FWD\_"}, and \texttt{"PT\_REV\_"} are all distinct.
In particular, $\mathtt{replica\_seed}(b, r, s) \ne \mathtt{swap\_seed}(b, s, r)$
for all $(b, r, s)$ because the prefixes differ even when the numeric arguments
coincide.
A test asserts that each prefix constant in source equals its expected string
and that $\mathtt{replica\_seed}(0, 0, 0)$ and $\mathtt{swap\_seed}(0, 0, 0)$
equal hard-coded expected values, preventing silent seed-compatibility drift.
\end{remark}

\subsection{Formal Pseudocode}

Algorithm~\ref{alg:pt} gives the complete \pt{} procedure.

\begin{algorithm}[t]
\caption{\pt{}($G$, $k$, $b$, $N$, $\varepsilon_c$, $\varepsilon_h$,
  $\mathtt{swap\_interval}$, $T$, $p$)}
\label{alg:pt}
\begin{algorithmic}[1]
\State \Comment{Geometric tolerance ladder}
\For{$i = 0$ \textbf{to} $N-1$}
  \State $\varepsilon_i \gets \varepsilon_c \times (\varepsilon_h / \varepsilon_c)^{i/(N-1)}$
\EndFor
\State \Comment{Initialise all replicas from the same METIS plan}
\State $\mathbf{a}_{\mathrm{init}} \gets \metis\text{-}k\text{-}{\rm plan}(G, k, \varepsilon_c, b)$
\For{$i = 0$ \textbf{to} $N-1$}
  \State $\mathtt{chain}_i \gets \fr\text{-}\mathrm{Chain}(G, k, \mathbf{a}_{\mathrm{init}},
    \varepsilon_i)$
\EndFor
\State $\mathcal{R} \gets \bigl[(\EC(\mathbf{a}_{\mathrm{init}}),\; 0,\;
  \mathbf{a}_{\mathrm{init}}.\mathrm{clone}())\bigr]$
\Comment{cold chain records: (EC, step, assignment)}
\For{$s = 1$ \textbf{to} $T$}
  \State \Comment{Step all replicas independently}
  \For{$i = 0$ \textbf{to} $N-1$}
    \State $r_s \gets \mathtt{replica\_seed}(b, i, s)$
    \State $\mathtt{rng\_fwd} \gets \mathrm{SmallRng}(\mathrm{SHA\text{-}256}(
      \texttt{"PT\_FWD\_"} \| r_s))$
    \State $\mathtt{rng\_rev} \gets \mathrm{SmallRng}(\mathrm{SHA\text{-}256}(
      \texttt{"PT\_REV\_"} \| r_s))$
    \State $\mathtt{chain}_i.\mathrm{step}(\mathtt{rng\_fwd},\; \mathtt{rng\_rev})$
  \EndFor
  \State $\mathcal{R}.\mathrm{push}\bigl((\EC(\mathtt{chain}_0.\mathbf{a}),\; s,\;
    \mathtt{chain}_0.\mathbf{a}.\mathrm{clone}())\bigr)$
  \State \Comment{Swap proposals every swap\_interval steps}
  \If{$s \bmod \mathtt{swap\_interval} = 0$}
    \For{$i = 0$ \textbf{to} $N-2$}
      \State $\mathtt{rng}_{\mathrm{swap}} \gets
        \mathrm{SmallRng}(\mathtt{swap\_seed}(b, s, i))$
      \State $\EC_i \gets \EC(\mathtt{chain}_i.\mathbf{a})$
      \State $\EC_j \gets \EC(\mathtt{chain}_{i+1}.\mathbf{a})$
      \State $\beta_i \gets 1/\varepsilon_i$; \quad $\beta_j \gets 1/\varepsilon_{i+1}$
      \State $\ell \gets (\beta_i - \beta_j) \times (\EC_i - \EC_j)$
      \If{$\mathtt{rng}_{\mathrm{swap}}.\mathrm{gen::<f64>}().\ln() < \min(0, \ell)$}
        \State \textbf{swap} $\mathtt{chain}_i.\mathbf{a} \leftrightarrow
          \mathtt{chain}_{i+1}.\mathbf{a}$
      \EndIf
    \EndFor
  \EndIf
\EndFor
\State \textbf{sort} $\mathcal{R}$ by $(\EC \;\mathrm{ASC},\; s \;\mathrm{ASC})$
\State \textbf{return} $\mathcal{R}[\lfloor p \times |\mathcal{R}| \rfloor].\mathbf{a}$
\end{algorithmic}
\end{algorithm}

\paragraph{Three design choices.}

\textbf{1. Geometric tolerance ladder.}
Equal spacing in $\ln\varepsilon$ gives approximately equal swap acceptance
rates across all adjacent pairs.
A ladder with equal absolute spacing $\Delta\varepsilon$ would concentrate swaps
near the cold end (where the tolerance difference matters most) and produce
near-universal acceptance near the hot end.
The geometric ladder distributes swap events more uniformly.

\textbf{2. Swap interval $= 10$.}
The swap interval balances two competing objectives.
If swaps are proposed too frequently (interval $= 1$), the chains do not have
time to diffuse between swap events; consecutive swaps on the same pair are
anti-correlated, reducing effective exploration.
If swaps are proposed too infrequently (interval $= 100$), the hot chain may
return to the cold chain's basin before the swap is proposed, losing the
diversity benefit.
Interval $= 10$ is the standard practical choice from the tempering
literature~\citep{earl2005} and is confirmed by the swap acceptance rates in
Section~\ref{sec:empirical}.

\textbf{3. Cold chain recording.}
The cold chain records all plans it visits, not only the plans where an
internal \fr{} proposal was accepted.
Regardless of whether the \fr{} proposal was accepted, the cold chain is in
some plan (either the proposed plan or the previous one); we record that plan.
This means $|\mathcal{R}| = T + 1$ (the initial plan plus one record per step).
The percentile selection at line~25 is over all cold chain states, including
states that persisted across multiple steps due to proposal rejection.

\paragraph{Cold chain acceptance rate vs.\ swap acceptance rate.}
Two acceptance rates are tracked and reported in the audit log:
\begin{itemize}
\item \texttt{swap\_acceptance\_rate}: Fraction of swap proposals that were
  accepted across all adjacent pairs. Measures chain exchange health.
\item \texttt{cold\_chain\_acceptance\_rate}: Fraction of \fr{} steps at replica
  0 where the internal proposal was accepted. Measures cold chain health.
\end{itemize}
A very low cold chain acceptance rate with a normal swap rate indicates that
the cold tolerance is too tight for this state: the cold chain cannot move
at $\varepsilon_{\mathrm{cold}}$, but the PT mechanism has no way to help
because the cold chain rejects proposals before swaps can inject diversity.
A normal cold chain acceptance rate with very low swap rate indicates that the
tolerance ladder is over-spaced: adjacent chains are too different in tolerance
to exchange efficiently.

\subsection{CLI Invocation and YAML Configuration}

\pt{} is exposed via the \texttt{--search parallel-tempering} flag:
\begin{verbatim}
redist state --state TX --year 2020 \
  --search parallel-tempering \
  --pt-replicas 4 \
  --pt-swap-interval 10 \
  --pt-cold-tol 0.005 \
  --pt-hot-tol 0.05 \
  --seeds 1000 \
  --percentile 0.0
\end{verbatim}

YAML configuration:
\begin{verbatim}
algorithm:
  structure: ratio-optimal
  weights: geographic
  search: parallel-tempering
  pt_replicas: 4
  pt_swap_interval: 10
  pt_cold_tolerance: 0.005
  pt_hot_tolerance: 0.05
  percentile: 0.0
\end{verbatim}

\subsection{Audit Chain}

Every run appends the following fields to \texttt{runs/\{label\}/\{year\}/index.json}:
\begin{verbatim}
"search": "parallel-tempering",
"n_replicas": 4,
"swap_interval": 10,
"cold_tolerance": 0.005,
"hot_tolerance": 0.05,
"percentile": 0.0,
"base_seed": 12345678,
"steps": 1000,
"swap_acceptance_rate": 0.21,
"cold_chain_acceptance_rate": 0.44,
"selected_step": 847,
"replica_seed_formula": "SHA-256('PT_REPLICA_' || r:u32le || '_' || s:u64le || '_' || b:u64le)",
"swap_seed_formula":    "SHA-256('PT_SWAP_'    || i:u32le || '_' || s:u64le || '_' || b:u64le)"
\end{verbatim}

The field \texttt{selected\_step} is the \emph{rank} of the returned plan in
the cold chain's EC-sorted record list ($0 =$ minimum EC), not a temporal step
index.
An auditor who re-runs the chain with the same \texttt{base\_seed} collects all
cold chain records, sorts by (EC ASC, step ASC), and confirms that the plan at
rank \texttt{selected\_step} matches the submitted plan.

\subsection{Computational Cost}

\pt{} costs $N_{\mathrm{replicas}} \times T$ \fr{} steps plus
$\lfloor T / \mathtt{swap\_interval} \rfloor \times (N_{\mathrm{replicas}}-1)$
swap evaluations.
Each swap evaluation requires computing $\EC$ for two chain states, which costs
$O(|E|)$ per plan.
In practice, $\EC$ can be maintained incrementally across \fr{} steps so that
the per-swap cost is $O(1)$ (only the boundary edges of the changed district
pairs need updating).
The dominant cost is the $N_{\mathrm{replicas}} \times T$ \fr{} steps, which is
$N$ times the cost of a single-chain run.
For $N=4$, the wall-clock cost is approximately $4\times$ the baseline single-chain
run time (assuming serial chain execution); with parallel execution via Rayon,
the wall-clock cost approaches the baseline (open question, noted in
Section~\ref{sec:conclusion}).
